% ======================================================================
%  RELATED WORK
% ======================================================================
\section{Related Work}\label{sec:related}

We organize related work along two axes: (1)~the signal generation method (rule-based z-score thresholds vs learned models) and (2)~the market structure (single-exchange different-asset pairs vs cross-exchange same-asset spreads). Our work occupies the intersection of \emph{learned spread forecasting} and \emph{cross-exchange same-asset trading} with \emph{minute-scale microstructure features}, a combination that, to the best of our knowledge, no prior study addresses.

\subsection{Rule-Based Z-Score Pairs Trading in Crypto}

All existing crypto pairs-trading studies use rolling z-score thresholds or cointegration-derived signals at 1--5 minute frequencies on single exchanges, trading \emph{different} assets whose co-movement may break down as fundamentals diverge~\cite{fil2020pairs, tadi2021evaluation, ko2023pairs, tadi2025copula, palazzi2025trading}. Our work differs in three structural ways: we trade the \emph{same} asset across \emph{different} exchanges (where the Law of One Price guarantees mean-reversion), we train a model to \emph{predict the future z-score} rather than using it as a mechanical threshold signal, and we incorporate L2~orderbook, trade flow, and funding rate features that cannot be backfilled from exchange APIs.

\subsection{ML Models for Spread Prediction}

A smaller body of work trains models to predict spread dynamics. Tsoku and Makatjane~\cite{tsoku2026dl} train a DNN+LSTM ensemble to forecast spread dynamics on cointegrated cryptocurrency pairs but operate at lower frequency, on a single exchange, without microstructure features. Fischer et al.~\cite{fischer2019statistical} train a random forest on minute-binned lagged returns of 40~crypto coins, finding that alpha falls from 20.5~bps to 3.8~bps under a one-minute execution delay---a decay profile that makes execution latency a first-order threat to validity at our collector cadence (Section~\ref{sec:limitations}). Liou et al.~\cite{liou2024hft} use LOB features for spread prediction in equities pairs trading, paralleling our orderbook imbalance inputs but on a single exchange. Han and Li~\cite{han2024lstm} use LSTM as a trade filter on CSI~300 futures, a framing analogous to our confidence threshold ($|\hat{z}| \geq 0.5$).

\subsection{ML-Enhanced Pairs Trading on Equities}

Several equity-market papers establish the dual-metric reporting norms we follow. Sarmento et al.~\cite{sarmento2024ml} find that model abstention improves profit-per-trade on 1-minute Brazilian stock ratios, paralleling our confidence filter. Shen et al.~\cite{shen2022arbitrage} achieve R$^2$ of 92\% on CSI~300 spot-futures spreads with LSTM, though the high R$^2$ reflects slow-moving equity index spreads; our noisy crypto spreads yield R$^2$ of 0.35 on filtered entries. Perrone et al.~\cite{perrone2026pairs} reformulate statistical arbitrage as a panel-level prediction task on daily equity data. These studies operate at lower frequencies, on single exchanges, and without non-backfillable microstructure features.

\subsection{Positioning Summary}

Table~\ref{tab:positioning} summarizes the key dimensions along which our work differs from the closest related papers.

\begin{table}[h]
\caption{Positioning of this work relative to closest related papers. Columns indicate whether each study uses a learned model (vs rule-based), operates cross-exchange (vs single-venue), uses L2 orderbook or microstructure features, evaluates on live (non-backtest) data, and reports both forecast and trading metrics.}
\label{tab:positioning}
\small
\begin{tabular}{lccccc}
\toprule
\textbf{Paper} & \textbf{Learned} & \textbf{Cross-} & \textbf{L2/} & \textbf{Live} & \textbf{Dual} \\
 & \textbf{model} & \textbf{exch.} & \textbf{micro.} & \textbf{eval} & \textbf{metrics} \\
\midrule
Fil \& Kristoufek      & \texttimes & \texttimes & \texttimes & \texttimes & \texttimes \\
Tadi \& Kortchemski     & \texttimes & \texttimes & \texttimes & \texttimes & partial \\
Ko et al.               & \texttimes & \texttimes & \texttimes & \texttimes & \checkmark \\
Tadi \& Witzany         & \texttimes & \texttimes & \texttimes & \texttimes & \checkmark \\
Fischer et al.          & \checkmark & \texttimes & \texttimes & \texttimes & partial \\
Tsoku \& Makatjane      & \checkmark & \texttimes & \texttimes & \texttimes & partial \\
Liou et al.             & \checkmark & \texttimes & \checkmark & \texttimes & partial \\
Han \& Li               & \checkmark & \texttimes & \texttimes & \texttimes & \texttimes \\
\midrule
\textbf{This work}      & \checkmark & \checkmark & \checkmark & \checkmark & \checkmark \\
\bottomrule
\end{tabular}
\end{table}

No prior study combines all five properties. The closest papers in individual dimensions are Tsoku and Makatjane~\cite{tsoku2026dl} (learned crypto spread forecast, but lower frequency and single exchange), Liou et al.~\cite{liou2024hft} (LOB features for spread prediction, but equities and single exchange), and Fischer et al.~\cite{fischer2019statistical} (minute-scale ML on crypto, but cross-sectional rank prediction rather than pairwise spread forecasting). Our contribution is at the intersection: a learned z-score forecast on cross-exchange same-asset spreads, trained on non-backfillable minute-scale microstructure features, validated live with dual forecast and trading metrics.
